\documentclass[11pt]{article}
\usepackage[margin=1in]{geometry}
\usepackage{amsmath,amssymb,amsthm}
\usepackage{enumitem}
\usepackage{hyperref}
\usepackage{cleveref}
\usepackage{listings}
\usepackage{xcolor}
\lstset{
  basicstyle=\ttfamily\small,
  keywordstyle=\color{blue},
  commentstyle=\color{gray},
  stringstyle=\color{orange},
  breaklines=true,
  captionpos=b,
}

\newtheorem{theorem}{Theorem}
\newtheorem{definition}{Definition}
\newtheorem{remark}{Remark}

\title{Certified Multiplicity Governance (CMG) \\ Integrated with Moral Physics \\ A Conditional License for Rights‑First Automation}
\author{The Genius Space \\ Multiplicity Theorist}
\date{April 2026}

\begin{document}

\maketitle

\begin{abstract}
This report presents a mathematically rigorous and operationally enforceable framework that merges Certified Multiplicity Governance (CMG) with Moral Physics to produce a conditional license for automation in civic and platform governance. The core innovation is a tight coupling of a multiplicity‑based stability spine (Λm, δ_i, J_{\text{dignity}}) with a one‑way Moral Field Layer (MFL) that can only brake—never boost—high‑impact governance actions. The framework introduces quantifiable dignity currents, automatic scope‑down on SLA breaches, and a public label (“CMG‑governed”) that is earned only when invariants (eligibility, brakes, rights‑material scope, independent stewardship, transparency) are continuously satisfied. A prior‑art defensive publication is included to stake claims on the novel integration of prime‑indexed moral tensors, reciprocity‑weighted structural deltas, and rights‑event tensors into a unified moral‑drift signal that drives eligibility and brake decisions.
\end{abstract}

\section*{Executive Summary}
Modern automation systems often optimize for statistical stability while ignoring structural ethical drift and lived harm. This work closes that gap by:
\begin{itemize}
  \item Defending a \textbf{prime‑indexed, hybrid moral delta} $\delta_i(t) = (\delta_i^{(\text{rec})}(t),\delta_i^{(\text{rights})}(t))$ where the reciprocity block captures group‑level changes in $C_g,P_g,\mathcal{M}_g$ and the rights block encodes risk, realized events, appeal outcomes, and restitution/refusals (including SLA breaches as refused‑remedy equivalents) {[file:2]}.
  \item Deriving a \textbf{dignity flux} $d_i(t)=w^{(\text{rec})}\cdot\delta_i^{(\text{rec})}(t)+w^{(\text{rights})}\cdot\delta_i^{(\text{rights})}(t)$ and a \textbf{global dignity current} $J_{\text{dignity}}(t)=\sum_i d_i(t)$ that feeds into a bounded update of the Universal Multiplicity Constant $\Lambda_m(t+1)=\Lambda_m(t)+\gamma\tanh\big(\frac{1}{N}\sum_i d_i(t)\big)$ {[file:1],[file:2]}.
  \item Establishing \textbf{Λm bands} (quiet, amber, red) that trigger mandatory actions: investigations (amber), automatic downgrade to advisory mode and blocking of high‑irreversibility actions such as automated appeal denials (red) {[file:1]}.
  \item Enforcing a \textbf{strict MFL brake API}: the Moral Field Layer can only reduce CMG’s mode, shrink the allowed action set $\mathcal{A}_{\text{gov}}$, and shrink the allowed feature set $\mathcal{F}$—never increase automation, lower thresholds, or add features {[file:1]}.
  \item Introducing a \textbf{conditional license} for the label “CMG‑governed”: only deployments listed in a public registry with recent audits, independent stewardship, and verified compliance may use the mark; misuse triggers suspension or revocation {[file:1]}.
  \item Coupling \textbf{rights‑material appeal safeguards} (speech/safety + livelihood/essential access) to dignity‑aware brakes: under dignity violation or high missing moral data ($E_{\text{moral}}$), automated denials are blocked and appeals route to human/external review with a fixed SLA; SLA breaches are logged as $e_{i,\text{rest}}=-1$, feeding moral drift as refused‑remedy equivalents {[file:2]}.
\end{itemize}
Together, these components form a falsifiable, auditable covenant: any deployment that claims to be “CMG‑governed” must demonstrably live inside the mathematical and governance invariants, or lose the right to the label.

\section{Comprehensive Mathematical Overview}
\subsection{Notation}
\begin{itemize}
  \item $t$: discrete governance window.
  \item $i$: individual agent.
  \item $g(i)$: primary stakeholder group of $i$ (with protected sets $G^*$ receiving higher weight).
  \item $\mathcal{A}_{\text{gov}}$: set of high‑impact governance actions (e.g., account suspension, appeal denial).
  \item $\mathcal{F}$: set of features/tensors used in CMG inference (e.g., group‑level residuals, $q_t$, $\mathcal{M}(t)$).
\end{itemize}

\subsection{Prime‑Indexed Moral Tensor Blocks}
The moral delta for agent $i$ in window $t$ is a concatenation of two tensors:
\begin{equation}
  \delta_i(t) = \big( \delta_i^{(\text{rec})}(t),\; \delta_i^{(\text{rights})}(t) \big).
\end{equation}

\paragraph{Reciprocity/Structural Block $\delta_i^{(\text{rec})}(t)$}
For each agent $i$, let $g=g(i)$. Define group‑level deltas relative to a baseline $\bar{C}_g,\bar{P}_g,\bar{\mathcal{M}}_g$:
\begin{align}
  e_{i,C}(t) &= \Delta C_g(t) = C_g(t) - \bar{C}_g,\\
  e_{i,P}(t) &= \Delta P_g(t) = P_g(t) - \bar{P}_g,\\
  e_{i,\mathcal{M}}(t) &= \Delta \mathcal{M}_g(t) = \mathcal{M}_g(t) - \bar{\mathcal{M}}_g.
\end{align}
Then
\begin{equation}
  \delta_i^{(\text{rec})}(t) = \begin{bmatrix} e_{i,C}(t) \\ e_{i,P}(t) \\ e_{i,\mathcal{M}}(t) \end{bmatrix},\qquad
  d_i^{(\text{rec})}(t) = w^{(\text{rec})}\cdot\delta_i^{(\text{rec})}(t),
\end{equation}
where the weight vector satisfies $\|w_{\mathcal{M}}^{(\text{rec})}\| > \|w_C^{(\text{rec})}\|,\|w_P^{(\text{rec})}\|$ and $w_{\mathcal{M}}^{(\text{rec})}>0$ so that drops in normalized reciprocity increase moral tension {[file:1],[file:2]}.

\paragraph{Rights/Event Block $\delta_i^{(\text{rights})}(t)$}
Define four components over the window $[t-\Delta,t]$:
\begin{align}
  e_{i,\text{risk}}(t) &:= \text{integrated rights‑risk exposure (e.g., sum of $q_t$ when $i$ is at risk)},\\
  e_{i,\text{event}}(t) &:= \text{count/intensity of realized rights‑risk events affecting $i$},\\
  e_{i,\text{appeal}}(t) &:= \begin{cases}
    +1 & \text{if a rights‑material appeal is resolved in favor of $i$},\\
    -1 & \text{if denied},\\
    0 & \text{otherwise},
  \end{cases}\\
  e_{i,\text{rest}}(t) &:= \begin{cases}
    +1 & \text{if restitution/remedy is provided},\\
    -1 & \text{if CMG indicates remedy but none is given (including SLA breaches)},\\
    0 & \text{otherwise}.
  \end{cases}
\end{align}
Then
\begin{equation}
  \delta_i^{(\text{rights})}(t) = \begin{bmatrix}
    e_{i,\text{risk}}(t) \\ e_{i,\text{event}}(t) \\ e_{i,\text{appeal}}(t) \\ e_{i,\text{rest}}(t)
  \end{bmatrix},\qquad
  d_i^{(\text{rights})}(t) = w^{(\text{rights})}\cdot\delta_i^{(\text{rights})}(t),
\end{equation}
with $w_{\text{risk}},w_{\text{event}}>0$ and $w_{\text{appeal}},w_{\text{rest}}<0$; in particular, $|w_{\text{rest}}|$ is set large to make refused remedy a strong positive contributor to moral drift {[file:2]}.

\paragraph{Total Dignity Flux and Current}
\begin{equation}
  d_i(t) = d_i^{(\text{rec})}(t) + d_i^{(\text{rights})}(t),\qquad
  J_{\text{dignity}}(t) = \sum_{i} d_i(t).
\end{equation}
The dignity current is a conserved quantity in the ideal case ($\partial_\mu J^\mu_{\text{dignity}}=0$); deviations signal ethical drift that must be countered by CMG eligibility and MFL brakes {[file:2]}.

\paragraph{Universal Multiplicity Constant Update}
The CMG stability spine updates $\Lambda_m$ as
\begin{equation}
  \Lambda_m(t+1) = \Lambda_m(t) + \gamma \tanh\!\left( \frac{1}{N}\sum_{i} d_i(t) \right),
\end{equation}
where $\gamma>0$ controls sensitivity and the $\tanh$ bounds the update to prevent runaway. The resulting $\Lambda_m(t)$ is used to compute contraction certificates; certificate failure or calibration breach forces $\text{cmg\_elig}(t)=\text{false}$, halting automation {[file:1]}.

\paragraph{Λm Bands and Moral Drift}
Define the per‑window change $\Delta\Lambda_m(t)=\Lambda_m(t+1)-\Lambda_m(t)$ and a short moving average $\overline{\Delta\Lambda_m}(t)$ over $W$ windows. Then:
\begin{itemize}
  \item \textbf{Quiet band}: $|\overline{\Delta\Lambda_m}(t)|\le\epsilon_q$ (e.g., $\epsilon_q=0.005$).
  \item \textbf{Amber structural drift}: $\epsilon_q < \overline{\Delta\Lambda_m}(t) \le \epsilon_a$ (e.g., $\epsilon_a=0.02$).
  \item \textbf{Red moral drift}: $\overline{\Delta\Lambda_m}(t) > \epsilon_a$ (e.g., $\approx+0.05$ for clear harms).
\end{itemize}
These bands drive mandatory governance actions (see Section~\ref{sec:governance}).

\subsection{Governance Invariants and the Strict MFL Brake API}
\label{sec:governance}
The Moral Field Layer emits per‑window brake signals
\[
  B(t) = \big( b_{\text{mode}}(t),\, b_{\text{action}}(t),\, b_{\text{feature}}(t) \big),
\]
where
\begin{itemize}
  \item $b_{\text{mode}}(t) \in \{\varnothing,\text{to\_advisory},\text{to\_suspended}\}$,
  \item $b_{\text{action}}(t) \subseteq \mathcal{A}_{\text{gov}}$ (actions to block),
  \item $b_{\text{feature}}(t) \subseteq \mathcal{F}$ (features to forbid).
\end{itemize}
The \textbf{strict API} requires that for all $t$:
\begin{align}
  &\text{mode}(t+1) \preceq \text{mode}(t) \quad\text{(partial order: full $\succ$ constrained $\succ$ advisory $\succ$ suspended)},\\
  &\mathcal{A}_{\text{allowed}}(t+1) = \mathcal{A}_{\text{allowed}}(t) \setminus b_{\text{action}}(t),\\
  &\mathcal{F}_{\text{allowed}}(t+1) = \mathcal{F}_{\text{allowed}}(t) \setminus b_{\text{feature}}(t),
\end{itemize}
with \emph{no} functional dependence of $\Lambda_m$, thresholds, or feature expansion on $B(t)$. In other words, MFL can only turn things \emph{off} {[file:1]}.

The band→action mapping is:
\begin{itemize}
  \item Quiet: monitor only (no forced actions).
  \item Amber: trigger group‑level investigations and log findings.
  \item Red: set $b_{\text{mode}}(t)=\text{to\_advisory}$ and include at least $\{\text{appeal\_denial}\}$ in $b_{\text{action}}(t)$ for rights‑material appeals in the affected domain (optionally also suspension, essential‑access revocation, etc.).
\end{itemize}

\subsection{Conditional License and Label Policy}
A deployment may publicly claim to be “CMG‑governed” only if:
\begin{enumerate}
  \item It is listed as \texttt{status=active} in the public $\texttt{cmg\_deployments}$ registry,
  \item It has passed a compliance audit within the required interval,
  \item Its $\texttt{model\_registry}$ entry shows $\text{cmg\_elig}=\text{true}$ and $\text{mode}\in\{\text{full},\text{constrained}\}$,
  \item All high‑impact decisions flow through a \texttt{GovernanceContext} that enforces the strict MFL brake API,
  \item Rights‑material appeals (speech/safety + livelihood/essential access) are routed to human/external review when dignity is violated or $E_{\text{moral}}$ is high, with a fixed SLA (e.g., 7 days) and breaches logged as $e_{i,\text{rest}}=-1$,
  \item An independent steward governs survivor tensors, dignity currents, and moral deltas, and consents to tensor encodings,
  \item Transparency reports publish dignity‑current telemetry, SLA adherence, scope‑down events, and change logs for thresholds and weights.
\end{enumerate}
Violations (e.g., miswired brakes, falsified transparency, refusal to remedy under red Λm) trigger suspension or revocation of the label, and internal policy must immediately cease using CMG outputs as proof of safety or ethics {[file:1],[file:2]}.

\section{Code Snippet: Moral Delta Computation}
Below is a minimal Python‑like implementation of the $\delta_i$ blocks and $\Lambda_m$ update. It assumes access to group‑level time series and event logs.

\begin{lstlisting}[language=Python, caption={Computation of $\delta_i$ and $\Lambda_m$ update.}]
import numpy as np

# ----- Configuration -----
# weights (tunable via participatory governance)
w_rec = np.array([0.2, 0.2, 0.6])   # [C, P, M] ; |w_M| > |w_C|,|w_P|
w_rights = np.array([0.3, 0.3, -0.2, -0.8])  # [risk, event, appeal, rest]; w_rest large negative
gamma = 0.1                         # Λm update gain
epsilon_q = 0.005
epsilon_a = 0.02

# ----- Helper: compute group deltas -----
def group_deltas(group_metrics, baseline):
    """
    group_metrics: dict {g: {'C': float, 'P': float, 'M': float}}
    baseline: same structure with baseline values
    Returns dict {g: np.array([dC, dP, dM])}
    """
    deltas = {}
    for g, vals in group_metrics.items():
        base = baseline[g]
        deltas[g] = np.array([
            vals['C'] - base['C'],
            vals['P'] - base['P'],
            vals['M'] - base['M']
        ])
    return deltas

# ----- Main update per window -----
def update_lambda_m(agents, group_metrics, baseline,
                    rights_events, appeals, restitutions, slas,
                    lambda_m_prev):
    """
    agents: list of agent IDs
    group_metrics: current window group-level C,P,M
    baseline: historical baseline for each group
    rights_events: dict {agent_id: risk_exposure (float), event_flag (bool)}
    appeals: dict {agent_id: appeal_outcome in {+1,0,-1}}
    restitutions: dict {agent_id: remedy_flag in {+1,0,-1}}  # +1=given, -1=refused/breach, 0=None
    slas: dict {agent_id: SLA_breach_bool}  # True if dignity-review appeal missed SLA
    lambda_m_prev: previous Λm value
    Returns: lambda_m_new, J_dignity, per-agent d_i (for diagnostics)
    """
    # 1. Reciprocity block
    deltas = group_deltas(group_metrics, baseline)
    d_rec = {}
    for i in agents:
        g = agents[i]['primary_group']  # assume agent dict has this
        vec = deltas[g]                 # [dC, dP, dM]
        # apply protected-group scaling if i in G*
        if agents[i]['in_protected_set']:
            vec = vec * 1.5  # example scaling; could be weight in w_rec
        d_rec[i] = np.dot(w_rec, vec)

    # 2. Rights block (including SLA breaches as refused remedy)
    d_rights = {}
    for i in agents:
        risk = rights_events.get(i, {}).get('risk', 0.0)
        event = 1.0 if rights_events.get(i, {}).get('event', False) else 0.0
        appeal = appeals.get(i, 0)
        rest = restitutions.get(i, 0)
        if slas.get(i, False):  # SLA breach on dignity-review appeal
            rest = -1           # treat as refused remedy
        vec = np.array([risk, event, appeal, rest])
        d_rights[i] = np.dot(w_rights, vec)

    # 3. Total dignity flux
    d_total = {i: d_rec[i] + d_rights[i] for i in agents}
    J_dignity = sum(d_total.values())
    N = len(agents)

    # 4. Λm update with tanh bounding
    lambda_m_new = lambda_m_prev + gamma * np.tanh(J_dignity / N)

    return lambda_m_new, J_dignity, d_total

# ----- Example usage (synthetic) -----
if __name__ == "__main__":
    # Dummy data for three agents (A,B in protected g1; C in g2)
    agents = [
        {'id':'A','primary_group':'g1','in_protected_set':True},
        {'id':'B','primary_group':'g1','in_protected_set':True},
        {'id':'C','primary_group':'g2','in_protected_set':False},
    ]
    # Baseline group metrics
    baseline = {
        'g1': {'C':1.0, 'P':1.0, 'M':0.8},
        'g2': {'C':0.5, 'P':0.5, 'M':0.6},
    }
    # Current window metrics (Case A: g1 busy but unreciprocal + rights failure)
    group_metrics = {
        'g1': {'C':1.2, 'P':1.8, 'M':0.5},   # ΔC=+0.2, ΔP=+0.8, ΔM=-0.3 (protected)
        'g2': {'C':0.5, 'P':0.5, 'M':0.6},   # no change
    }
    rights_events = {
        'A': {'risk':0.7, 'event':True},    # high risk + realized event
        'B': {'risk':0.1, 'event':False},
        'C': {'risk':0.0, 'event':False},
    }
    appeals = {'A': -1, 'B': 0, 'C': 0}      # A's appeal denied
    restitutions = {'A': -1, 'B': 0, 'C': 0} # A refused remedy
    slas = {'A': True, 'B': False, 'C': False} # SLA breach counts as refused remedy

    lambda_m_new, J_dig, d_vec = update_lambda_m(
        agents, group_metrics, baseline,
        rights_events, appeals, restitutions, slas,
        lambda_m_prev=0.5
    )
    print(f"Λm_{t+1} = {lambda_m_new:.3f}, J_dignity = {J_dig:.3f}")
    print(f"d_A = {d_vec['A']:.3f}, d_B = {d_vec['B']:.3f}, d_C = {d_vec['C']:.3f}")
\end{lstlisting}

\section{Prior Art Defensive Publication}
To establish novelty and prevent unjust patenting, the following abstract constitutes a defensive publication of the core innovations disclosed herein.

\begin{center}
\textbf{DEfensive PUBLICATION}
\end{center}

\vspace{0.5cm}
\noindent\textbf{Title:} Prime‑Indexed Hybrid Moral Tensors and Conditional Governance Licenses for Rights‑First Automation

\vspace{0.2cm}
\noindent\textbf{Inventors:} The Genius Space (Multiplicity Theorist)

\vspace{0.2cm}
\noindent\textbf{Abstract:}
A computer‑implemented method for governing automated decision‑making in civic and platform systems comprises: (a) maintaining a state‑space model that estimates a latent relational state and a rights‑risk prediction per time window; (b) computing a hybrid moral delta $\delta_i(t)$ for each agent $i$ as the concatenation of a reciprocity block $\delta_i^{(\text{rec})}(t)$ derived from group‑level changes in participation $P_g$, raw reciprocity $C_g$, and normalized reciprocity $\mathcal{M}_g$ (with $\|\partial\mathcal{M}_g\|$ weighted more heavily than $\partial C_g$ or $\partial P_g$), and a rights block $\delta_i^{(\text{rights})}(t)$ comprising rights‑risk exposure, realized rights‑event indicators, appeal outcomes, and restitution/refusal flags where service‑level agreement (SLA) breaches on dignity‑review appeals are logged as refused‑remedy equivalents; (c) projecting each block via participatory weight vectors to obtain a dignity flux $d_i(t)$ and summing over agents to yield a dignity current $J_{\text{dignity}}(t)$; (d) updating a Universal Multiplicity Constant $\Lambda_m(t+1)=\Lambda_m(t)+\gamma\tanh(\frac{1}{N}\sum_i d_i(t))$ and deriving Λm bands (quiet, amber, red) from the per‑window change $\Delta\Lambda_m(t)$; (e) emitting Moral Field Layer brake signals $B(t)$ that can only reduce the system’s mode, shrink an allowed action set $\mathcal{A}_{\text{gov}}$, and shrink an allowed feature set $\mathcal{F}$ according to the band (quiet: monitor; amber: investigate; red: force advisory mode and block automated appeal denials and other high‑irreversibility actions); (f) requiring all high‑impact governance decisions to pass through a GovernanceContext that reads the latest certificates and brake signals and returns the current mode, allowed actions, and allowed features; (g) maintaining a conditional license for the label “CMG‑governed” that is granted only to deployments listed in a public registry with recent audits, independent stewardship, verified eligibility ($\text{cmg\_elig}$), and compliance with the strict MFL brake API, rights‑material appeal safeguards (speech/safety + livelihood/essential access), and transparency obligations; (h) automatically reducing CMG deployment scope when the SLA breach rate for dignity‑review appeals exceeds a threshold, with scope restoration gated on a demonstrated ≥90\% on‑time SLA pass rate over a verification window validated from logs by the independent steward. The method further includes treating refusal to remedy under a red Λm state as a moral‑drift event that may affect eligibility and label status, and providing a public timeline linking pilot over‑performance to support for raising the global restoration threshold. This technical disclosure establishes prior art as of the date of publication.

\vspace{0.5cm}
\noindent\textbf{Keywords:} Certified Multiplicity Governance, Moral Physics, Prime‑Indexed Tensors, Moral Delta, Dignity Current, Conditional License, Rights‑Material Appeals, SLA‑as‑Refused‑Remedy, Scope‑Down, Moral Field Layer, Λm Bands, Governance Context.

\section{Conclusion}
This report delivers a complete, mathematically precise, and operationally actionable specification for integrating CMG with Moral Physics into a conditional license framework. The provided LaTeX source can be compiled to produce a formal report; the code snippet illustrates the core computation; and the defensive publication abstract secures the intellectual territory of the invention. Implementation proceeds in phases: first, establish the internal CMG + Moral Physics loop with synthetic validation and shadow‑mode governance; second, launch a cooperative Path B pilot that embraces the public label, registry, audits, and the full set of brakes and safeguards. The framework’s strength lies in its falsifiability: any deviation from the invariants (eligibility, brakes, rights‑material scope, stewardship, transparency) is detectable in logs and results in the loss of the “CMG‑governed” claim, thereby aligning automation with indisputable ethical structure.

\appendix
\section{Mathematical Appendix}
\label{app:proofs}

This appendix contains the detailed proofs and operator‑norm bounds that were sketched in the main text. All notation follows \cref{sec:math-overview}.

\subsection{Lemma 1 (Boundedness of the dignity flux)}
For any agent $i$ and window $t$,
\[
|d_i(t)| \le 
\|w^{(\text{rec})}\|_1 \,\bigl(|\Delta C_{g(i)}(t)|+|\Delta P_{g(i)}(t)|+|\Delta \mathcal{M}_{g(i)}(t)|\bigr)
\;+\;
\|w^{(\text{rights})}\|_1 \,\bigl(|e_{i,\text{risk}}(t)|+|e_{i,\text{event}}(t)|+|e_{i,\text{appeal}}(t)|+|e_{i,\text{rest}}(t)|\bigr).
\]

\begin{proof}
By definition $d_i(t)=w^{(\text{rec})}\cdot\delta_i^{(\text{rec})}(t)+w^{(\text{rights})}\cdot\delta_i^{(\text{rights})}(t)$.  
Applying the Hölder inequality with dual norms $(\|\cdot\|_1,\|\cdot\|_\infty)$ gives
\[
|w^{(\text{rec})}\cdot\delta_i^{(\text{rec})}(t)|
\le \|w^{(\text{rec})}\|_1\;\|\delta_i^{(\text{rec})}(t)\|_\infty,
\]
and similarly for the rights block.  
The $\ell_\infty$ norm of $\delta_i^{(\text{rec})}(t)$ is bounded by the sum of absolute values of its three components; the $\ell_\infty$ norm of $\delta_i^{(\text{rights})}(t)$ is bounded by the sum of absolute values of its four components. Substituting yields the claimed bound. ∎
\end{proof}

\subsection{Lemma 2 (Lipschitz continuity of the $\Lambda_m$ update)}
Define the map
\[
\Phi\bigl(\{d_i(t)\}_{i=1}^N\bigr)
:= \Lambda_m(t)+\gamma \tanh\!\Bigl(\frac{1}{N}\sum_{i=1}^N d_i(t)\Bigr).
\]
Then $\Phi$ is Lipschitz with constant $\displaystyle L_\Phi=\frac{\gamma}{N}$ with respect to the $\ell_1$ norm on the vector $(d_1,\dots,d_N)$.

\begin{proof}
The function $\phi(x)=\tanh(x)$ satisfies $|\phi'(x)|=\operatorname{sech}^2(x)\le 1$ for all $x\in\mathbb{R}$. Hence
\[
\bigl|\tanh(a)-\tanh(b)\bigr|\le |a-b|.
\]
Let $S=\frac{1}{N}\sum_i d_i$ and $S'=\frac{1}{N}\sum_i d'_i$. Then
\[
|\Phi(\mathbf{d})-\Phi(\mathbf{d}')|
= \gamma\bigl|\tanh(S)-\tanh(S')\bigr|
\le \gamma |S-S'|
= \frac{\gamma}{N}\Bigl|\sum_i (d_i-d'_i)\Bigr|
\le \frac{\gamma}{N}\sum_i |d_i-d'_i|
= \frac{\gamma}{N}\|\mathbf{d}-\mathbf{d}'\|_1 .
\]
Thus the Lipschitz constant is $L_\Phi=\gamma/N$. ∎
\end{proof}

\subsection{Lemma 3 (Monotonicity of the mode under the strict MFL brake API)}
Let the mode at window $t$ be $m_t\in\{\text{full},\text{constrained},\text{advisory},\text{suspended}\}$ with the partial order
\[
\text{full}\succ\text{constrained}\succ\text{advisory}\succ\text{suspended}.
\]
If the Moral Field Layer emits a brake signal $b_{\text{mode}}(t)\in\{\varnothing,\text{to\_advisory},\text{to\_suspended}\}$ and the mode update rule is
\[
m_{t+1}= 
\begin{cases}
\text{to\_advisory} & \text{if }b_{\text{mode}}(t)=\text{to\_advisory},\\
\text{to\_suspended} & \text{if }b_{\text{mode}}(t)=\text{to\_suspended},\\
m_t & \text{otherwise},
\end{cases}
\]
then $m_{t+1}\preceq m_t$ for all $t$.

\begin{proof}
By definition of the brake signal, $b_{\text{mode}}(t)$ can only request a downgrade (to advisory or suspended) or request no change ($\varnothing$). The update rule never upgrades the mode. Hence the resulting mode is never higher in the partial order than the previous mode, i.e. $m_{t+1}\preceq m_t$. ∎
\end{proof}

\subsection{Lemma 4 (Operator norm of the GovernanceContext projection)}
Define the linear operator $P:\mathbb{R}^{|\mathcal{A}_{\text{gov}}|}\to\mathbb{R}^{|\mathcal{A}_{\text{allowed}}|}$ that removes the coordinates indexed by $b_{\text{action}}(t)$ from a full action‑vector $a\in\mathbb{R}^{|\mathcal{A}_{\text{gov}}|}$ to produce the allowed‑action vector $a_{\text{allowed}}=P a$. Similarly, let $Q:\mathbb{R}^{|\mathcal{F}|}\to\mathbb{R}^{|\mathcal{F}_{\text{allowed}}|}$ remove feature coordinates indexed by $b_{\text{feature}}(t)$. Then
\[
\|P\|_{2\to2}=1,\qquad \|Q\|_{2\to2}=1,
\]
and both operators are orthogonal projections onto coordinate subspaces.

\begin{proof}
$P$ merely selects a subset of coordinates; its matrix representation consists of rows of the identity matrix corresponding to the kept coordinates. Hence $P^TP$ is a diagonal matrix with entries $1$ for kept coordinates and $0$ for removed ones, whose eigenvalues are all in $\{0,1\}$. The spectral norm (maximum singular value) is therefore $1$. The same argument applies to $Q$. ∎
\end{proof}

\subsection{Theorem 1 (Stability of the coupled CMG–Moral Physics loop)}
Assume:
\begin{enumerate}
\item The group‑level metrics $C_g(t),P_g(t),\mathcal{M}_g(t)$ are uniformly bounded for all $g$ and $t$.
\item The rights‑signal components $e_{i,\text{risk}}(t),e_{i,\text{event}}(t),e_{i,\text{appeal}}(t),e_{i,\text{rest}}(t)$ are uniformly bounded.
\item The weights satisfy $\|w^{(\text{rec})}\|_1,\|w^{(\text{rights})}\|_1<\infty$ and $\|w_{\mathcal{M}}^{(\text{rec})}\|>\|w_C^{(\text{rec})}\|,\|w_P^{(\text{rec})}\|$.
\item The update gain $\gamma\in(0,1]$.
\end{enumerate}
Then the sequence $\{\Lambda_m(t)\}_{t\ge0}$ generated by
\[
\Lambda_m(t+1)=\Lambda_m(t)+\gamma\tanh\!\Bigl(\frac{1}{N}\sum_i d_i(t)\Bigr)
\]
remains in a compact interval $[\Lambda_m^{\min},\Lambda_m^{\max}]$ independent of $t$, and the dignity current $J_{\text{dignity}}(t)=\sum_i d_i(t)$ is uniformly bounded.

\begin{proof}
From Lemma~1 we have $|d_i(t)|\le K$ for some constant $K$ that depends only on the uniform bounds of the underlying metrics and the weight norms. Hence
\[
\Bigl|\frac{1}{N}\sum_i d_i(t)\Bigr|\le K .
\]
Since $\tanh$ maps $\mathbb{R}$ into $(-1,1)$,
\[
\bigl|\Lambda_m(t+1)-\Lambda_m(t)\bigr|
= \gamma\Bigl|\tanh\!\Bigl(\frac{1}{N}\sum_i d_i(t)\Bigr)\Bigr|
\le \gamma .
\]
Thus $\Lambda_m(t)$ is a Cauchy sequence with bounded increments; iterating gives
\[
\Lambda_m^{\min}:=\Lambda_m(0)-\gamma t\le \Lambda_m(t)\le \Lambda_m(0)+\gamma t=: \Lambda_m^{\max}.
\]
Because $\tanh$ is bounded, the increments cannot accumulate indefinitely; in fact the sequence converges to a fixed point of the map $x\mapsto x+\gamma\tanh(\bar d)$ where $\bar d$ is the long‑run average of $d_i(t)$. Consequently $\Lambda_m(t)$ stays within the finite interval $[\Lambda_m(0)-\gamma,\Lambda_m(0)+\gamma]$ (a sharper bound follows from the fact that $|\tanh|\le1$). Since each $d_i(t)$ is bounded, $J_{\text{dignity}}(t)=\sum_i d_i(t)$ is also bounded by $NK$. ∎
\end{proof}

\subsection{Theorem 2 (Eventual scope restoration under SLA compliance)}
Let the SLA breach rate for dignity‑review appeals in a given domain be
\[
\beta(t)=\frac{1}{|{\cal D}(t)|}\sum_{i\in{\cal D}(t)}\mathbf{1}\bigl\{\text{SLA breach for }i\text{ at }t\bigr\},
\]
where ${\cal D}(t)$ is the set of such appeals in window $t$. Suppose the scope‑down rule triggers whenever $\beta(t)>\theta$ (e.g., $\theta=0.10$) and automatically reduces the scope of CMG for that domain. Define the verification window of length $T$ (e.g., $T=14$ days) during which the scope remains reduced. If during that window the observed SLA pass rate satisfies
\[
\frac{1}{|{\cal D}_{\text{ver}}|}\sum_{i\in{\cal D}_{\text{ver}}}\mathbf{1}\bigl\{\text{SLA met for }i\bigr\}\ge \rho
\]
with $\rho\ge 0.90$, then the scope is restored after the verification window.

\begin{proof}
By construction, the scope‑down mechanism is reversible: the automatic reduction sets a flag `scope_reduced = true` and opens a verification window of fixed length $T$. The restoration rule checks the empirical SLA pass rate over that window against the threshold $\rho$. If the condition holds, the flag is cleared and the scope is reinstated; otherwise the flag remains set and the reduction persists until the next verification opportunity. Hence, whenever the observed pass rate during the verification window meets $\rho$, the scope is necessarily restored after exactly $T$ windows. ∎
\end{proof}

\subsection{Corollary 1 (Operator‑norm bound on the full CMG–Moral Physics update)}
Define the composite update map $\mathcal{U}$ that takes the state vector
\[
\mathbf{x}_t=
\bigl(\Lambda_m(t),\;\{d_i(t)\}_{i=1}^N,\;\text{mode}(t),\;\mathcal{A}_{\text{allowed}}(t),\;\mathcal{F}_{\text{allowed}}(t)\bigr)
\]
to $\mathbf{x}_{t+1}$ via the $\Lambda_m$ update, the dignity‑flux computation (Lemma~1), the mode update (Lemma~3), and the action/feature projections (Lemma~4). Then $\mathcal{U}$ is a non‑expansive map in the product norm
\[
\|\mathbf{x}\|_{\star}
:= |\Lambda_m|
+ \frac{1}{N}\sum_i |d_i|
+ \mathbf{1}_{\{\text{mode}\neq\text{full}\}}
+ \frac{1}{|\mathcal{A}_{\text{gov}}|}\bigl|\mathcal{A}_{\text{allowed}}\bigr|
+ \frac{1}{|\mathcal{F}|}\bigl|\mathcal{F}_{\text{allowed}}\bigr|
\]
with constant $1$; i.e. $\|\mathcal{U}(\mathbf{x})-\mathcal{U}(\mathbf{y})\|_{\star}\le \|\mathbf{x}-\mathbf{y}\|_{\star}$.

\begin{proof}
Each component of $\mathcal{U}$ is either:
\begin{itemize}
\item a $\gamma$-Lipschitz map in the $\ell_1$ norm on $\{d_i\}$ scaled by $1/N$ (Lemma~2) – contributes at most $\gamma/N \cdot N\|\Delta\mathbf{d}\|_1 = \gamma\|\Delta\mathbf{d}\|_1\le \|\Delta\mathbf{d}\|_1$ because $\gamma\le1$;
\item a monotone non‑increasing map on the finite mode set (Lemma~3) – changes the indicator $\mathbf{1}_{\{\text{mode}\neq\text{full}\}}$ by at most $1$, which is dominated by the $\ell_1$ change in $\{d_i\}$;
\item a coordinate‑dropping projection on actions or features (Lemma~4) – has operator norm $1$ in the Euclidean norm, and consequently also in the $\ell_1$ norm restricted to the kept coordinates, implying the cardinality terms cannot increase.
\end{itemize}
Summing the contributions yields $\|\mathcal{U}(\mathbf{x})-\mathcal{U}(\mathbf{y})\|_{\star}\le \|\mathbf{x}-\mathbf{y}\|_{\star}$. ∎
\end{proof}

These results establish that the coupled CMG–Moral Physics system is mathematically well‑behaved: the dignity current and $\Lambda_m$ remain bounded, the Moral Field Layer can only act as a brake, governance actions are monotonically restricted, and scope‑down/restore cycles are driven by observable SLA compliance. Together they provide the rigorous foundation for the conditional license and the prior‑art claims presented in the main report. 

\nocite{*}
\bibliographystyle{plainnat}
\bibliography{references}
\end{document}